% !TEX root =./ECE444_Practice_main.tex
%
% L26 practice set (Beam Squint and Quantization) -- grating lobes, the spacing
% criterion, beam squint across a band, and B-bit phase quantization on the
% 8-element ADALM-PHASER row.

\fullwidth{\textbf{Learning Objective 3.8: } Predict and quantify the three ways a real
steered array departs from the ideal pattern: grating lobes from element spacing, beam
squint from operating away from the phase-set frequency, and the pointing and sidelobe
limits of a $B$-bit phase shifter.
\\[4pt]\normalfont\small Show your work. A \textbf{Documentation} line at the end
is required (write \textbf{None} if you did not collaborate).}

\begin{questions}

%%%%%%%%%%%%%%%%%%%%%%%%%%%%%%%%%%%%%%%%%%%%%%%%%%%%%%%%%%%%%%%%%%%%%%%%%%%%%%
\question \textbf{Where the array factor repeats.} The ADALM-PHASER carries 8 patch
elements on a $14\ \text{mm}$ pitch. In a thinning demonstration only every third element
is fed, so the fed elements sit $42\ \text{mm}$ apart. Work at the array design center,
$10.3\ \ghz$, where $\lambda = 29.1\ \text{mm}$.

\begin{parts}

	\part In one or two sentences, explain why a grating lobe reaches the same height as
	the main lobe, while a sidelobe does not.
		\begin{solution}[1.2in]
		At a grating-lobe angle the path-length difference between neighboring elements is
		a whole wavelength rather than zero, so every element contribution arrives in phase
		exactly as it does at the commanded angle. The array cannot distinguish the two
		directions, and the summation reaches its maximum value $N$ in both. A sidelobe is a
		partial cancellation of that same summation, so it is always below the peak.
		\end{solution}

	\begin{minipage}[t]{0.58\linewidth}
		\part With the beam at broadside, find every grating-lobe angle in visible space.
		\begin{solution}[1.3in]
		\eq{ \sin\theta_g &= \sin\theta_0 \pm m\frac{\lambda}{d} = 0 \pm m\frac{29.1}{42} \\
		     &= \pm 0.693 m }
		For $m=1$, $\theta_g = \arcsin(\pm 0.693) = \pm 43.9\mydeg$. For $m=2$ the required
		sine is $\pm 1.386$, which is not a real angle, so the pattern has three full-height
		beams and no more.
		\end{solution}
	\end{minipage}\hfill%
	\begin{minipage}[t]{0.37\linewidth}
		\raggedleft \vspace{12pt}
		$\theta_g$ = \ansbox[1.5in]{$\pm 43.9\mydeg$}
	\end{minipage}

	\begin{minipage}[t]{0.58\linewidth}
		\part Now steer the beam to $\theta_0 = +20\mydeg$. Find every grating lobe that
		remains in visible space.
		\begin{solution}[1.4in]
		\eq{ \sin\theta_g &= \sinp{20\mydeg} \pm 0.693 = 0.342 \pm 0.693 }
		The plus sign gives $1.035$, which is outside $[-1,1]$, so that lobe has left visible
		space. The minus sign gives $-0.351$, so
		\eq{ \theta_g = \arcsin(-0.351) = -20.5\mydeg }
		Steering the main beam to $+20\mydeg$ pushed one grating lobe past the horizon and
		pulled the other in from $-43.9\mydeg$ to $-20.5\mydeg$. Grating lobes move
		uniformly in $\sin\theta$, not in degrees.
		\end{solution}
	\end{minipage}\hfill%
	\begin{minipage}[t]{0.37\linewidth}
		\raggedleft \vspace{12pt}
		$\theta_g$ = \ansbox[1.5in]{$-20.5\mydeg$}
	\end{minipage}

	\part A classmate proposes applying a Hann taper to remove the extra beams. Explain in
	one or two sentences why that does not work.
		\begin{solution}[1.2in]
		A taper controls the sidelobe skirts that surround a single beam by smoothing the
		aperture illumination. A grating lobe is a second main beam produced by the element
		spacing, and it is present with any amplitude distribution, including a uniform one.
		The taper lowers the sidelobes around both beams equally and leaves both peaks at
		full height. The only cure is closer element spacing.
		\end{solution}

\end{parts}

\newpage
%%%%%%%%%%%%%%%%%%%%%%%%%%%%%%%%%%%%%%%%%%%%%%%%%%%%%%%%%%%%%%%%%%%%%%%%%%%%%%
\question \textbf{Choosing the element spacing.} An X-band linear array is being laid out.
Use $c = 3\times10^8\ \m/\text{s}$ throughout.

\begin{parts}

	\part State the criterion that keeps grating lobes out of visible space over a scan
	range $\pm\theta_0$, and explain why half-wavelength spacing is the usual default.
		\begin{solution}[1.3in]
		The criterion is
		\eq{ d < \frac{\lambda}{1 + \left|\sin\theta_0\right|} }
		Taking the scan range to its limit, $\theta_0 \rightarrow 90\mydeg$, gives
		$d < \lambda/2$. A half-wavelength lattice therefore has no grating lobe at any
		commandable scan angle, so it is the spacing that removes the question entirely
		rather than tying the array to one scan sector.
		\end{solution}

	\begin{minipage}[t]{0.58\linewidth}
		\part The array must scan to $\pm 60\mydeg$ and operate up to $10.5\ \ghz$. Find the
		largest usable element spacing in millimeters.
		\begin{solution}[1.3in]
		Use the highest frequency, where $\lambda$ is shortest and the criterion is tightest:
		\eq{ \lambda = \frac{3\times10^8}{10.5\times10^9} = 28.6\ \text{mm} }
		\eq{ d < \frac{28.6}{1 + \sinp{60\mydeg}} = \frac{28.6}{1.866} = 15.3\ \text{mm} }
		\end{solution}
	\end{minipage}\hfill%
	\begin{minipage}[t]{0.37\linewidth}
		\raggedleft \vspace{12pt}
		$d_{\max}$ = \ansbox[1.4in]{$15.3\ \text{mm}$}
	\end{minipage}

	\begin{minipage}[t]{0.58\linewidth}
		\part A competing layout uses $d = 20\ \text{mm}$. At $10.5\ \ghz$, how far off
		broadside can that array scan before a grating lobe enters visible space?
		\begin{solution}[1.3in]
		\eq{ \frac{\lambda}{d} = \frac{28.6}{20} = 1.429 }
		\eq{ 1 + \left|\sin\theta_0\right| < 1.429 \quad \Rightarrow \quad
		     \left|\sin\theta_0\right| < 0.429 }
		\eq{ \theta_0 < \arcsin(0.429) = 25.4\mydeg }
		\end{solution}
	\end{minipage}\hfill%
	\begin{minipage}[t]{0.37\linewidth}
		\raggedleft \vspace{12pt}
		$\theta_{0,\max}$ = \ansbox[1.4in]{$25.4\mydeg$}
	\end{minipage}

	\part The PHASER uses $d = 14\ \text{mm}$. State its scan limit from the same criterion,
	and give one sentence on what the design gives up by choosing a spacing this tight.
		\begin{solution}[1.3in]
		At $10.5\ \ghz$, $\lambda/d = 28.6/14 = 2.04$, so the criterion requires
		$1 + \left|\sin\theta_0\right| < 2.04$, which every real angle satisfies. The array
		has no grating lobe at any scan angle in its band.

		What it gives up is aperture. With the element count fixed at 8, tighter spacing
		means a shorter array, $L = Nd = 112\ \text{mm}$ instead of $160\ \text{mm}$ at
		$20\ \text{mm}$ pitch, so the beam is wider and the directivity lower. For a fixed
		element count, wide scan coverage and a narrow beam pull in opposite directions.
		\end{solution}

\end{parts}

\newpage
%%%%%%%%%%%%%%%%%%%%%%%%%%%%%%%%%%%%%%%%%%%%%%%%%%%%%%%%%%%%%%%%%%%%%%%%%%%%%%
\question \textbf{Beam squint across a band.} An 8-element array with $d = 14\ \text{mm}$
operates from $10.0$ to $10.5\ \ghz$. The beam-steering computer sets the element phases
for $\theta_0 = 45\mydeg$ at the band center, $f_0 = 10.25\ \ghz$. Use
$c = 3\times10^8\ \m/\text{s}$.

\begin{parts}

	\part Explain in one or two sentences why a broadside beam does not squint at all.
		\begin{solution}[1.1in]
		At broadside the commanded phase ramp is zero, because every element is equidistant
		from a source on boresight. Squint arises when a fixed phase ramp is interpreted at a
		different wavenumber, and scaling a ramp of zero still gives zero. The condition
		$\sin\theta = (f_0/f)\sin\theta_0$ returns $\theta = 0$ for any frequency.
		\end{solution}

	\begin{minipage}[t]{0.58\linewidth}
		\part Find the beam angle at the low band edge, $10.0\ \ghz$.
		\begin{solution}[1.3in]
		\eq{ \sin\theta &= \frac{f_0}{f}\sinp{45\mydeg}
		     = \frac{10.25}{10.00}(0.7071) \\
		     &= 1.025(0.7071) = 0.7248 }
		\eq{ \theta = \arcsin(0.7248) = 46.5\mydeg }
		The beam has squinted $+1.5\mydeg$ toward the horizon.
		\end{solution}
	\end{minipage}\hfill%
	\begin{minipage}[t]{0.37\linewidth}
		\raggedleft \vspace{12pt}
		$\theta$ = \ansbox[1.4in]{$46.5\mydeg$}
	\end{minipage}

	\begin{minipage}[t]{0.58\linewidth}
		\part Find the beam angle at the high band edge, $10.5\ \ghz$, and state the total
		angular walk of the beam across the band.
		\begin{solution}[1.4in]
		\eq{ \sin\theta &= \frac{10.25}{10.50}(0.7071) = 0.9762(0.7071) = 0.6903 }
		\eq{ \theta = \arcsin(0.6903) = 43.6\mydeg }
		The total walk from one band edge to the other is
		\eq{ 46.5\mydeg - 43.6\mydeg = 2.9\mydeg }
		\end{solution}
	\end{minipage}\hfill%
	\begin{minipage}[t]{0.37\linewidth}
		\raggedleft \vspace{12pt}
		walk = \ansbox[1.4in]{$2.9\mydeg$}
	\end{minipage}

	\begin{minipage}[t]{0.58\linewidth}
		\part The half-power beamwidth of this array at $45\mydeg$ is
		$\theta_\text{HP} \approx 0.886\lambda/(Nd\cosp{\theta_0})$. Express the walk from
		part (c) as a fraction of a beamwidth, and recommend a fix if the requirement were
		to hold the walk under 2\% of a beamwidth.
		\begin{solution}[1.5in]
		At $10.25\ \ghz$, $\lambda = 29.3\ \text{mm}$ and $Nd = 112\ \text{mm}$:
		\eq{ \theta_\text{HP} &= \frac{0.886(29.3)}{112\cosp{45\mydeg}}
		     = \frac{25.9}{79.2} = 0.327\ \text{rad} = 18.8\mydeg }
		\eq{ \frac{2.9\mydeg}{18.8\mydeg} = 0.15 }
		The beam walks about 15\% of a beamwidth, which costs a few tenths of a decibel at
		the aimpoint and is usually acceptable for a narrowband link. Holding it under 2\%
		is a factor of seven, and no phase setting achieves it because the walk is fixed by
		the fractional bandwidth and $\tanp{\theta_0}$. The fix is to steer with true time
		delay instead of phase, at the element or at the subarray level, so that the delay
		compensates the path difference itself and every frequency in the band points the
		same way.
		\end{solution}
	\end{minipage}\hfill%
	\begin{minipage}[t]{0.37\linewidth}
		\raggedleft \vspace{12pt}
		walk / $\theta_\text{HP}$ = \ansbox[1.2in]{$0.15$}
	\end{minipage}

\end{parts}

\newpage
%%%%%%%%%%%%%%%%%%%%%%%%%%%%%%%%%%%%%%%%%%%%%%%%%%%%%%%%%%%%%%%%%%%%%%%%%%%%%%
\question \textbf{What the phase shifter can actually do.} The ADAR1000 beamformer on the
PHASER uses a $B$-bit phase shifter. Work at the HB100 source frequency,
$10.525\ \ghz$, where $\lambda = 28.5\ \text{mm}$ and $d/\lambda = 0.491$.

\begin{parts}

	\begin{minipage}[t]{0.58\linewidth}
		\part Find the phase step (LSB) for $B = 6$ and for the ADAR1000's $B = 7$.
		\begin{solution}[1.1in]
		\eq{ \text{LSB} = \frac{360\mydeg}{2^B} }
		\eq{ B = 6: \quad \frac{360\mydeg}{64} = 5.625\mydeg }
		\eq{ B = 7: \quad \frac{360\mydeg}{128} = 2.8125\mydeg }
		\end{solution}
	\end{minipage}\hfill%
	\begin{minipage}[t]{0.37\linewidth}
		\raggedleft \vspace{12pt}
		$B=6$: \ansbox[1.1in]{$5.625\mydeg$} \\[10pt]
		$B=7$: \ansbox[1.1in]{$2.8125\mydeg$}
	\end{minipage}

	\begin{minipage}[t]{0.58\linewidth}
		\part With $B = 7$, find the pointing step near broadside and at
		$\theta_0 = 60\mydeg$.
		\begin{solution}[1.5in]
		The inter-element phase is $\Delta\phi = 360\mydeg (d/\lambda)\sin\theta_0$, so a
		change of one LSB moves $\sin\theta_0$ by
		\eq{ \Delta(\sin\theta_0) = \frac{\text{LSB}}{360\mydeg (d/\lambda)}
		     = \frac{2.8125}{360(0.491)} = 0.0159 }
		Near broadside $\sin\theta_0 \approx \theta_0$ in radians, so
		\eq{ \Delta\theta_0 = 0.0159\ \text{rad} = 0.91\mydeg }
		Off broadside divide by $\cos\theta_0$:
		\eq{ \Delta\theta_0 = \frac{0.0159}{\cosp{60\mydeg}} = 0.0318\ \text{rad}
		     = 1.82\mydeg }
		Both are small next to the array's $13.2\mydeg$ beamwidth.
		\end{solution}
	\end{minipage}\hfill%
	\begin{minipage}[t]{0.37\linewidth}
		\raggedleft \vspace{12pt}
		broadside: \ansbox[1.1in]{$0.91\mydeg$} \\[10pt]
		at $60\mydeg$: \ansbox[1.1in]{$1.82\mydeg$}
	\end{minipage}

	\begin{minipage}[t]{0.58\linewidth}
		\part Use the rule of thumb $\text{QSLL} \approx -6B\ \db$ for $B = 3$ and $B = 7$,
		and state which sidelobe mechanism dominates in each case for a uniform 8-element
		array.
		\begin{solution}[1.4in]
		\eq{ B = 3: \quad \text{QSLL} \approx -18\ \db }
		\eq{ B = 7: \quad \text{QSLL} \approx -42\ \db }
		A uniform 8-element array has its own first sidelobe at about $-13\ \text{dB}$. At three
		bits the quantization lobes are already below that, and at seven bits they are far
		below it, so in both cases the ordinary sidelobe structure of the uniform
		illumination dominates. Quantization only takes over below about two bits, where the
		rule predicts $-12\ \db$ and the array's own sidelobes are no longer the largest
		feature.
		\end{solution}
	\end{minipage}\hfill%
	\begin{minipage}[t]{0.37\linewidth}
		\raggedleft \vspace{12pt}
		$B=3$: \ansbox[1.1in]{$-18\ \db$} \\[10pt]
		$B=7$: \ansbox[1.1in]{$-42\ \db$}
	\end{minipage}

	\part A classmate argues that because the phase shifter resolves only $2.8125\mydeg$, a
	pattern null on the PHASER can never be driven deeper than the $21\ \text{dB}$ a sweep
	shows. Explain in two or three sentences why that reasoning is wrong, and state what
	does set the measured depth.
		\begin{solution}[1.5in]
		Rounding the null weights onto the $2.8125\mydeg$ phase grid, with gain steps of about
		1\%, still produces a notch near $-48\ \text{dB}$ at the angle the weights were
		designed for. Quantization displaces the notch by a fraction of a degree from its
		commanded angle rather than filling it in, so it is a pointing error and not a depth
		limit.

		What caps the measurement is the receiver. The beam sweep's noise floor sits about
		$23\ \db$ below the uniform-taper peak, and the null weights give up roughly
		$2\ \db$ of main-lobe gain, leaving about $21\ \db$ of usable dynamic range. The
		pattern goes deeper than that; the sweep cannot follow it down.
		\end{solution}

\end{parts}

\newpage
%%%%%%%%%%%%%%%%%%%%%%%%%%%%%%%%%%%%%%%%%%%%%%%%%%%%%%%%%%%%%%%%%%%%%%%%%%%%%%
\question \textbf{Reading a sweep and sizing a design.} Each part below reports something
a student measured on the PHASER, or asks for a design decision. The array is 8 elements
on a $14\ \text{mm}$ pitch unless stated otherwise.

\begin{parts}

	\begin{minipage}[t]{0.58\linewidth}
		\part A broadside sweep against the HB100 at $10.525\ \ghz$
		($\lambda = 28.5\ \text{mm}$) shows three peaks of equal height, at $0\mydeg$ and
		$\pm 42\mydeg$. What effective element spacing does that indicate, and how were the
		elements fed?
		\begin{solution}[1.4in]
		\eq{ \sin\theta_g = \frac{\lambda}{d} \quad \Rightarrow \quad
		     d = \frac{\lambda}{\sinp{42\mydeg}} = \frac{28.5}{0.669} = 42.6\ \text{mm} }
		That is three times the $14\ \text{mm}$ pitch, so every third element was fed and the
		other five were switched off. Equal peak heights confirm grating lobes rather than
		sidelobes.
		\end{solution}
	\end{minipage}\hfill%
	\begin{minipage}[t]{0.37\linewidth}
		\raggedleft \vspace{12pt}
		$d_\text{eff}$ = \ansbox[1.4in]{$42.6\ \text{mm}$}
	\end{minipage}

	\begin{minipage}[t]{0.58\linewidth}
		\part The array is commanded to $+45\mydeg$ with phases set at $10.525\ \ghz$, but
		the sweep peaks at $+47.9\mydeg$ with the source running at $10.025\ \ghz$. Name the
		effect and confirm the measurement.
		\begin{solution}[1.3in]
		This is beam squint: the phase ramp was computed at one frequency and the signal
		arrived at another.
		\eq{ \sin\theta &= \frac{10.525}{10.025}\sinp{45\mydeg}
		     = 1.0499(0.7071) = 0.7424 }
		\eq{ \theta = \arcsin(0.7424) = 47.9\mydeg }
		The measurement matches the prediction, so nothing is wrong with the array.
		\end{solution}
	\end{minipage}\hfill%
	\begin{minipage}[t]{0.37\linewidth}
		\raggedleft \vspace{12pt}
		$\theta$ = \ansbox[1.4in]{$47.9\mydeg$}
	\end{minipage}

	\part At $B = 2$ the sweep shows a lobe near $-55\mydeg$ about $7\ \db$ below the peak,
	while the rule of thumb predicts $-12\ \text{dB}$. Explain the discrepancy in one or two
	sentences.
		\begin{solution}[1.2in]
		The $-6B$ figure is an RMS estimate derived for a large aperture, where the staircase
		phase error repeats over many periods and the resulting lobes average out to the
		predicted level. An 8-element array holds only a couple of periods of that sawtooth,
		so individual lobes scatter several decibels either side of the estimate. The rule
		gives the expected level, not a bound on any single lobe.
		\end{solution}

	\begin{minipage}[t]{0.58\linewidth}
		\part A new array must scan to $\pm 50\mydeg$ over $10.0$ to $10.5\ \ghz$ and hold
		every sidelobe below $-25\ \text{dB}$. State the maximum element spacing and the minimum
		number of phase bits, and say which of the three defects the bandwidth requirement
		drives.
		\begin{solution}[1.6in]
		Spacing, evaluated at the top of the band where $\lambda = 28.6\ \text{mm}$:
		\eq{ d < \frac{28.6}{1 + \sinp{50\mydeg}} = \frac{28.6}{1.766} = 16.2\ \text{mm} }
		Bits, from the quantization sidelobe rule:
		\eq{ -6B < -25 \quad \Rightarrow \quad B > 4.2 \quad \Rightarrow \quad B = 5 }
		The ADAR1000's seven bits clear this with margin. The bandwidth requirement drives
		beam squint: the band is 4.9\% wide, so at $50\mydeg$ the beam walks roughly
		$(0.024)\tanp{50\mydeg} = 0.029\ \text{rad} = 1.6\mydeg$ either side of center. That
		is a small fraction of the beamwidth here, so phase steering is adequate; a wider
		band or a wider scan would force true time delay.
		\end{solution}
	\end{minipage}\hfill%
	\begin{minipage}[t]{0.37\linewidth}
		\raggedleft \vspace{12pt}
		$d_{\max}$ = \ansbox[1.2in]{$16.2\ \text{mm}$} \\[10pt]
		$B_{\min}$ = \ansbox[1.2in]{$5$ bits}
	\end{minipage}

\end{parts}

\vspace{1.5em}
\noindent\textbf{Documentation:} \rule[-2pt]{0.6\linewidth}{0.4pt}

\end{questions}
